\section{Implementation}
\label{sec:implementation}
\subsection{Edge and Cloud Runtime}
The primary deployment runs the student visual tower and feature codec on
a Jetson AGX Orin and the $\pi_{0.5}$ policy on an A40 cloud server.
The robot-side pipeline prepares each camera input, executes the student
and encoder, and serializes the entropy-coded payload. The server
reconstructs visual tokens and supplies them through the frozen policy's
visual-input interface. The integration uses asynchronous LeRobot
inference; it changes the visual transport path rather than introducing
a new action scheduler.

In the recorded SO-101 component benchmark, two camera images are resized
from $640\times480$ to $224\times168$ before the codec. This common resize
is applied to every transport mode and is outside the reported codec
timing boundary. The student and codec execute through a CUDA Graph in
the deployed edge path. The compiled cloud $\pi_{0.5}$ implementation uses
fixed buffers for visual-token injection.

\subsection{Entropy Coding and Numeric Precision}
The entropy implementation uses static integer tables at both endpoints.
The deployment tests compare encoded symbols and byte streams across the
edge and receiver, separating lossless coding from lossy representation
learning. Floating-point precision remains part of the codec
configuration: our diagnostics found that changing encoder
precision between training and deployment could change quantization
decisions. Precision must therefore be fixed together with the
checkpoint rather than treated as an interchangeable runtime option.

The component benchmark and deployed path use different synchronization
boundaries. Timing each component with an explicit device synchronization
includes overhead that is absent from a captured graph. We use the full
deployed-path measurement for comparisons of complete encoding and keep
the synchronized measurements as component diagnostics. They are not
summed to estimate deployed latency.

\subsection{Measurement Boundaries}
The latency replay records edge encoding, cloud decoding, total cloud
processing, and complete request latency. Its network column is computed
as the remainder after subtracting edge and cloud processing from the
request duration; it includes transfer, response return, and
serialization effects. It is not an independent one-way network probe.

For the replay, the first five of 200 uploads per method are discarded as
warm-up, leaving 195 measurements. Quantiles are computed per upload.
The measurements characterize a recorded visual sequence through the
actual client and server. They do not represent 195 independent physical
task trials. The replay table uses KiB because its generating script
divides payload bytes by 1,024; simulation payload summaries are labeled
separately to retain their original accounting.
